\begin{frame}[fragile]{OpenMDAO Group class}
  \begin{lstlisting}[basicstyle=\ttfamily\tiny]
  # Groups assemble and connect components together,
  # and have to deal with the flow of data between components

  class QuuxGroup(Group):

      def initialize(self):
          self.options.declare("foo", default=10.0, types=float)

      def setup(self):
          self.add_input("bar",  val=1.0, units="m")

          self.add_subsystem(
              "Foo",
              FooComponent(foo),
              promotes_inputs  = [("bar", "foo_bar")],
              promotes_outputs = [("xyz", "foo_xyz")],
          )

          self.add_subsystem(
              "Faz",
              FazComponent(foo),
              promotes_inputs  = [("bar", "faz_bar")],
              promotes_outputs = [("xyz", "faz_xyz")],
          )

          self.connect("Foo.ijk", "Faz.baz")
  \end{lstlisting}
\end{frame}

\begin{frame}[fragile]{OpenMDAO Group class}
  \begin{lstlisting}
  # Groups can promote inputs and outputs of their
  # components, in which case the promoted
  # input (output) becomes the input (ouptut)
  # of the whole group viewed as a component

  self.add_subsystem(
      "Foo",
      FooComponent(foo),
      promotes_inputs  = [("bar", "foo_bar")],
      promotes_outputs = [("xyz", "foo_xyz")],
      # . . . . . . . . . ^~~~~~~~~~~~~~~~~^
      # This tuple syntax instructs OpenMDAO
      # to rename inputs and outputs
      # to avoid naming clashes
  )
  \end{lstlisting}
\end{frame}

\begin{frame}[fragile]{OpenMDAO Group class}
  \begin{lstlisting}
  # Groups can connect inputs and outputs
  # of their components

  self.connect("Foo.ijk", "Faz.baz")

  # In the ideal case, inputs and outputs
  # are mapped cleanly without circular dependencies
  \end{lstlisting}
\end{frame}

\begin{frame}[fragile]{OpenMDAO Group class}
  \begin{lstlisting}
  # These connections can sometimes create
  # circular relations, which have to be
  # handled with nonlinear solvers

  self.nonlinear_solver
    = om.NewtonSolver(solve_subsystems=False)

  self.nonlinear_solver.options["maxiter"] = 30
  self.nonlinear_solver.options["rtol"]    = 1e-6
  self.nonlinear_solver.options["iprint"]  = 0

  self.nonlinear_solver.linesearch
    = om.ArmijoGoldsteinLS()

  self.linear_solver = om.DirectSolver()

  # This is an example of a very basic Newton's
  # method solver from our Propulsion group.
  # It is simple and efficient
  # under the right conditions
  \end{lstlisting}
\end{frame}
